\documentclass[aspectratio=169,10pt]{beamer}
\usetheme{Madrid}
\usecolortheme{seahorse}
\setbeamertemplate{footline}[frame number]

\usepackage{amsmath}
\usepackage{amssymb}
\usepackage{amsthm}
\usepackage{mathtools}
\usepackage{listings}
\usepackage{xcolor}
\usepackage{graphicx}

\graphicspath{{figures/}}

\title{Chapter 17: Differentiability of Convex Functions}
\subtitle{Subdifferential Examples \& Properties}
\author{Convex Analysis and Monotone Operator Theory}
\date{Pages 533--565}

% Define theorem styles
\newtheorem{prop}{Proposition}
\newtheorem{lem}{Lemma}
\newtheorem{cor}{Corollary}
\newtheorem{thm}{Theorem}
\newtheorem{defn}{Definition}
\newtheorem{exmpl}{Example}

% Listings style
\lstset{
  basicstyle=\ttfamily\scriptsize,
  keywordstyle=\color{blue},
  commentstyle=\color{green!50!black}\itshape,
  stringstyle=\color{red},
  breaklines=true,
  numbers=left,
  numberstyle=\tiny\color{gray},
  frame=single,
  language=Python
}

\begin{document}

\frame{\titlepage}

\section{Introduction}
\begin{frame}
  \frametitle{Chapter Overview}
  \begin{itemize}
    \item \textbf{Main Topic:} Differentiability and subdifferentiability of convex functions
    \item \textbf{Key Concepts:}
      \begin{itemize}
        \item Directional derivatives
        \item Gâteaux and Fréchet differentiability
        \item Subdifferentials and subgradients
        \item Characterizations of convexity
      \end{itemize}
    \item \textbf{Applications:}
      \begin{itemize}
        \item Optimization theory
        \item Monotone operator theory
        \item Numerical algorithms
      \end{itemize}
  \end{itemize}
\end{frame}

\section{17.1: Directional Derivatives}
\begin{frame}[fragile]
  \frametitle{Directional Derivatives (Section 17.1)}

  \begin{defn}[Directional Derivative]
    For $f: \mathcal{H} \to \mathbb{R} \cup \{+\infty\}$ proper and convex, with $x \in \text{dom} f$ and $y \in \mathcal{H}$:
    \[
    f'(x;y) = \inf_{\alpha \in \mathbb{R}_{++}} \frac{f(x + \alpha y) - f(x)}{\alpha}
    \]
  \end{defn}

  \begin{prop}[Proposition 17.2]
    Let $f: \mathcal{H} \to [-\infty, +\infty]$ be proper and convex. Then:
    \begin{enumerate}
      \item[(i)] $f'(x; y)$ exists in $[-\infty, +\infty]$
      \item[(ii)] $f'(x; y - x) + f(x) \leq f(y)$
      \item[(iii)] $\nabla f$ is sublinear if $f$ is differentiable
      \item[(iv)] $f'(x; 0) = 0$ and $f'(x; \cdot)$ is sublinear
    \end{enumerate}
  \end{prop}
\end{frame}

\begin{frame}
  \frametitle{Key Properties of Directional Derivatives}

  \begin{prop}[Proposition 17.3]
    Let $f: \mathcal{H} \to [-\infty, +\infty]$ be proper and convex, and $x \in \text{dom} f$. Then:
    \[
    x \in \text{Argmin} f \iff f'(x; \cdot) \geq 0
    \]
  \end{prop}

  \begin{prop}[Proposition 17.4]
    Let $f: \mathcal{H} \to [-\infty, +\infty]$ be convex and $x \in \mathcal{H}$. If $f$ is Gâteaux differentiable at $x$, then:
    \[
    x \in \text{Argmin} f \iff \nabla f(x) = 0
    \]
  \end{prop}

  \begin{itemize}
    \item The directional derivative characterizes whether a point is a minimum
    \item Gâteaux differentiability provides a linear characterization
  \end{itemize}
\end{frame}

\section{17.2: Characterizations of Convexity}
\begin{frame}
  \frametitle{Convexity Characterizations (Section 17.2)}

  \begin{prop}[Proposition 17.7]
    Let $f: \mathcal{H} \to [-\infty, +\infty]$ be proper with open convex domain. If $f$ is Gâteaux differentiable on $\text{dom} f$, then the following are equivalent:
    \begin{enumerate}
      \item[(i)] $f$ is convex
      \item[(ii)] $\langle y - x \mid \nabla f(x) \rangle + f(x) \leq f(y)$ for all $x, y$
      \item[(iii)] $\nabla f$ is monotone
    \end{enumerate}

    \vspace{0.5cm}
    If $f$ is twice Gâteaux differentiable, then equivalently:
    \begin{enumerate}
      \item[(iv)] $\langle z \mid \nabla^2 f(x) z \rangle \geq 0$ for all $z$
    \end{enumerate}
  \end{prop}
\end{frame}

\begin{frame}
  \frametitle{First-Order Characterization Example}

  \begin{exmpl}[Example 17.8]
    Let $A \in \mathcal{B}(\mathcal{H})$ and $f: \mathcal{H} \to \mathbb{R}: x \mapsto \langle x \mid Ax \rangle$. Then $f$ is convex if and only if:
    \[
    \langle x \mid (A + A^*) x \rangle \geq 0 \quad \forall x \in \mathcal{H}
    \]
  \end{exmpl}

  \begin{exmpl}[Example 17.9]
    Given $h: \mathcal{H} \to [-\infty, +\infty]$ Fréchet differentiable on open convex domain $D$, define:
    \[
    f(x) = \begin{cases}
      h(x), & x \in D \\
      \lim_{\alpha \downarrow 0} h((1-\alpha)y + \alpha x), & x \in \text{bdry } D \\
      +\infty, & x \notin \overline{D}
    \end{cases}
    \]
    Then $f \in \Gamma_0(\mathcal{H})$ and $f|_D = h|_D$ with lower semicontinuity.
  \end{exmpl}
\end{frame}

\section{17.3: Strict Convexity}
\begin{frame}
  \frametitle{Strict Convexity Characterizations (Section 17.3)}

  \begin{prop}[Proposition 17.10]
    Let $f: \mathcal{H} \to [-\infty, +\infty]$ be proper with open convex domain. If $f$ is Gâteaux differentiable on $\text{dom} f$, then these are equivalent:
    \begin{enumerate}
      \item[(i)] $f$ is strictly convex
      \item[(ii)] $x \neq y \Rightarrow \langle y - x \mid \nabla f(x) \rangle + f(x) < f(y)$
      \item[(iii)] $x \neq y \Rightarrow \langle x - y \mid \nabla f(x) - \nabla f(y) \rangle > 0$
    \end{enumerate}

    If $f$ is twice Gâteaux differentiable:
    \begin{enumerate}
      \item[(iv)] $z \neq 0 \Rightarrow \langle z \mid \nabla^2 f(x) z \rangle > 0$ for all $z \in \mathcal{H}$
    \end{enumerate}
  \end{prop}
\end{frame}

\begin{frame}
  \frametitle{Strict Convexity Examples}

  \begin{exmpl}[Example 17.11]
    Let $A \in \mathcal{B}(\mathcal{H})$ and $f: \mathcal{H} \to \mathbb{R}: x \mapsto \langle x \mid Ax \rangle$. Then $f$ is strictly convex if and only if $A + A^*$ is strictly monotone.
  \end{exmpl}

  \begin{exmpl}[Example 17.13]
    The function
    \[
    h: \mathbb{R}^2 \to [-\infty, +\infty]: (\xi, \eta) \mapsto \begin{cases}
      \eta^2 + \eta^2/\xi, & \xi > 0, \eta > 0 \\
      +\infty, & \text{otherwise}
    \end{cases}
    \]
    is strictly convex in the interior of its domain.
  \end{exmpl}
\end{frame}

\section{17.4: Directional Derivatives and Subgradients}
\begin{frame}[fragile]
  \frametitle{Definition and Basic Properties (Section 17.4)}

  \begin{defn}[Descent Direction]
    Let $f: \mathcal{H} \to [-\infty, +\infty]$ be proper and convex, with $x \in \text{dom} f$. Then $y \in \mathcal{H}$ is a descent direction of $f$ at $x$ if there exists $\varepsilon \in \mathbb{R}_{++}$ such that:
    \[
    \forall \alpha \in [0, \varepsilon] \quad f(x + \alpha y) < f(x)
    \]
  \end{defn}

  \begin{prop}[Proposition 17.21]
    Let $f: \mathcal{H} \to [-\infty, +\infty]$ be proper and convex. Then:
    \begin{enumerate}
      \item[(i)] $y$ is a descent direction at $x$ iff $f'(x; y) < 0$
      \item[(ii)] If $f$ is differentiable at $x \notin \text{Argmin} f$, then $-\nabla f(x)$ is a descent direction
    \end{enumerate}
  \end{prop}
\end{frame}

\begin{frame}
  \frametitle{Steepest Descent Direction}

  \begin{prop}[Proposition 17.22 -- Steepest Descent]
    Let $f: \mathcal{H} \to [-\infty, +\infty]$ be proper and convex with $x \in (\text{cont } f) \setminus (\text{Argmin } f)$. Set $u = P_{\partial f(x)}(0)$ and $z = -u/\|u\|$. Then $z$ is the unique minimizer of $f'(x; \cdot)$ over $B(0; 1)$, satisfying:
    \[
    f'(x; z) = \max_{z \in \partial f(x)} \langle z \mid \partial f(x) \rangle = -\|u\|
    \]
  \end{prop}

  \begin{itemize}
    \item The steepest descent direction is the negative of the projection onto the subdifferential
    \item This characterizes the direction of maximum decrease
  \end{itemize}
\end{frame}

\begin{frame}
  \frametitle{Chebyshev Center Example}

  \begin{prop}[Proposition 17.25 -- Chebyshev Center]
    Let $C$ be a nonempty compact convex subset of $\mathcal{H}$, and define:
    \[
    f: \mathcal{H} \to \mathbb{R}: x \mapsto \max_{y \in C} \|x - y\|^2
    \]

    Then:
    \begin{enumerate}
      \item[(i)] $f$ is continuous, strongly convex, and supercoercive
      \item[(ii)] $\text{dom } \partial f = \mathcal{H}$, $\text{gra } \partial f$ is closed, and the sets $(\partial f(x))_{x \in \mathcal{H}}$ are compact
      \item[(iv)] $\partial f(x) = 2(x - \text{conv } \partial_C(x))$
      \item[(v)] $f$ has a unique minimizer $r$, the Chebyshev center, characterized by $r \in \text{conv } \partial_C(r)$
    \end{enumerate}
  \end{prop}
\end{frame}

\section{17.5: Directional Derivatives and Convexity}
\begin{frame}
  \frametitle{Convexity from Directional Derivatives (Section 17.5)}

  \begin{prop}[Proposition 17.27]
    Let $f: \mathcal{H} \to [-\infty, +\infty]$ be proper with $x_0 \in \text{dom} f$ and $x_1 \in \mathcal{H}$ such that $[x_0, x_1] \subset \text{dom} f$. Suppose that $f|_{[x_0, x_1]}$ is lower semicontinuous and that for every $x \in [x_0, x_1]$, $f'(x; x_1 - x_0) \leq 0$. Then $f(x_1) \leq f(x_0)$.
  \end{prop}

  \begin{prop}[Proposition 17.29]
    Let $f: \mathcal{H} \to [-\infty, +\infty]$ be proper with $\text{dom} f$ convex and $f|_{\text{dom } f}$ lower semicontinuous. Suppose that for every $x, y \in \text{dom} f$:
    \[
    f'(x; y - x) \leq -f'(y; x - y)
    \]

    Then $f$ is convex.
  \end{prop}
\end{frame}

\begin{frame}
  \frametitle{Directional Derivatives and Convexity (continued)}

  \begin{cor}[Corollary 17.30]
    Let $f: \mathbb{R} \to [-\infty, +\infty]$ be proper with $I = \text{dom} f$ an interval. Suppose that $f|_I$ is lower semicontinuous and that for every $x, y \in I$ with $x < y$:
    \begin{enumerate}
      \item[(i)] $f_-'(x)$ and $f_-'(y)$ exist
      \item[(ii)] $f_-'(x) \leq f_-'(y)$ (monotonicity)
    \end{enumerate}

    Then $f$ is convex.
  \end{cor}

  This is the univariate analogue of Proposition 17.29.
\end{frame}

\section{17.6: Gâteaux and Fréchet Differentiability}
\begin{frame}
  \frametitle{Differentiability Concepts (Section 17.6)}

  \begin{defn}[Gâteaux and Fréchet Differentiability]
    Let $f: \mathcal{H} \to \mathbb{R} \cup \{+\infty\}$ be proper and convex at $x$:
    \begin{itemize}
      \item \textbf{Gâteaux differentiable:} $f'(x; y) = \langle y \mid \nabla f(x) \rangle$ for some $\nabla f(x) \in \mathcal{H}$
      \item \textbf{Fréchet differentiable:} $\lim_{0 \neq y \to 0} \frac{f(x+y) - f(x) - \langle y \mid \nabla f(x) \rangle}{\|y\|} = 0$
    \end{itemize}
  \end{defn}

  \begin{prop}[Proposition 17.31]
    Let $f: \mathcal{H} \to [-\infty, +\infty]$ be proper and convex, and $x \in \text{dom} f$. Then:
    \begin{enumerate}
      \item[(i)] If $f$ is Gâteaux differentiable at $x$, then $\partial f(x) = \{\nabla f(x)\}$
      \item[(ii)] If $x \in \text{cont } f$ and $\partial f(x)$ is a singleton $u$, then $f$ is Gâteaux differentiable and $\nabla f(x) = u$
    \end{enumerate}
  \end{prop}
\end{frame}

\begin{frame}
  \frametitle{Chain Rule for Differentiable Functions}

  \begin{prop}[Proposition 17.32]
    Let $f: \mathcal{H} \to \mathbb{R}$ be continuous and convex, and $\phi: \mathbb{R} \to \mathbb{R}$ be convex and differentiable with $\phi' \geq 0$. If $\phi \circ f$ is Gâteaux differentiable at $x$, then:
    \[
    \partial(\phi \circ f)(x) = \phi'(f(x)) \partial f(x)
    \]
  \end{prop}

  \begin{exmpl}[Example 17.33]
    Let $C$ be a closed convex set, $\phi: \mathbb{R} \to \mathbb{R}$ convex and differentiable on $\mathbb{R}_+$ with $\phi' > 0$ there. For $x \in \mathcal{H}$:
    \[
    \partial(\phi \circ d_C)(x) = \begin{cases}
      \{\phi'(d_C(x))(x - P_{C} x)\}, & x \notin C \\
      N_C \cap B(0; \phi'(0)), & x \in \text{bdry } C \\
      \{0\}, & x \in \text{int } C
    \end{cases}
    \]
  \end{exmpl}
\end{frame}

\begin{frame}
  \frametitle{Gâteaux vs Fréchet Differentiability}

  \begin{prop}[Proposition 17.39]
    Let $f \in \Gamma_0(\mathcal{H})$ with $x \in \text{int dom } f$. Then the following are equivalent:
    \begin{enumerate}
      \item[(i)] $f$ is Gâteaux differentiable at $x$
      \item[(ii)] Every selection of $\partial f$ is strong-to-weak continuous at $x$
      \item[(iii)] There exists a selection of $\partial f$ that is strong-to-weak continuous at $x$
    \end{enumerate}
  \end{prop}

  \begin{prop}[Proposition 17.41]
    Let $f \in \Gamma_0(\mathcal{H})$ with $x \in \text{int dom } f$. Then the following are equivalent:
    \begin{enumerate}
      \item[(i)] $f$ is Fréchet differentiable at $x$
      \item[(ii)] Every selection of $\partial f$ is continuous at $x$
      \item[(iii)] There exists a selection of $\partial f$ that is continuous at $x$
    \end{enumerate}
  \end{prop}
\end{frame}

\begin{frame}
  \frametitle{Finite-Dimensional Results}

  \begin{cor}[Corollary 17.44]
    If $\mathcal{H}$ is finite-dimensional, then Gâteaux and Fréchet differentiability are the same for functions in $\Gamma_0(\mathcal{H})$.
  \end{cor}

  \begin{prop}[Proposition 17.45]
    Assume $\mathcal{H}$ is finite-dimensional. Let $f \in \Gamma_0(\mathcal{H})$, $x \in \mathcal{H}$, and $u \in \mathcal{H}$. Then:
    \[
    \partial f(x) = \{u\} \iff f \text{ is Fréchet differentiable at } x \text{ and } u = \nabla f(x)
    \]
  \end{prop}

  \begin{itemize}
    \item In finite dimensions, the notions of differentiability coincide
    \item A singleton subdifferential characterizes Fréchet differentiability
  \end{itemize}
\end{frame}

\begin{frame}
  \frametitle{Monotone Operators Example}

  \begin{prop}[Proposition 17.36]
    Let $A \in \mathcal{B}(\mathcal{H})$ be monotone and self-adjoint. Suppose that $\text{ran } A$ is closed. Define:
    \[
    q_A: \mathcal{H} \to \mathbb{R}: x \mapsto \frac{1}{2} \langle x \mid Ax \rangle, \quad q_{A^\dagger}: \mathcal{H} \to \mathbb{R}: x \mapsto \frac{1}{2} \langle x \mid A^\dagger x \rangle
    \]

    Then:
    \begin{enumerate}
      \item[(i)] $q_A$ is convex, continuous, Fréchet differentiable, with $\nabla q_A = A$
      \item[(ii)] $q_{A^\dagger} = q_A \circ A^\dagger = \iota_{\ker A} \square q_{A^\dagger}$
      \item[(iii)] $q_{A^\dagger}^* = \iota_{\text{ran } A} + q_A^\dagger$
      \item[(iv)] $q_{A^\dagger}^* \circ A = q_A$
    \end{enumerate}
  \end{prop}
\end{frame}

\section{17.7: Differentiability and Continuity}
\begin{frame}
  \frametitle{Continuity from Differentiability (Section 17.7)}

  \begin{exmpl}[Example 17.47]
    Let $C$ be the closed cone of Example 8.42 and $f = \iota_C$. Then:
    \begin{itemize}
      \item $0 \in (\text{core } C) \setminus (\text{int } C)$
      \item $\partial f(0) = \{0\}$ is a singleton
      \item $f$ is Gâteaux differentiable at 0 with $\nabla f(0) = 0$
      \item $f$ is lower semicontinuous at 0
      \item But $0 \notin \text{int dom } f$, so $f$ is not continuous at 0
    \end{itemize}
  \end{exmpl}

  This shows that Gâteaux differentiability does not imply continuity.
\end{frame}

\begin{frame}
  \frametitle{Main Continuity Results}

  \begin{prop}[Proposition 17.48]
    Let $f: \mathcal{H} \to [-\infty, +\infty]$ be convex and Gâteaux differentiable at $x \in \text{dom} f$. Then:
    \begin{enumerate}
      \item[(i)] $f$ is lower semicontinuous at $x$
      \item[(ii)] If $\mathcal{H}$ is finite-dimensional, then $x \in \text{int dom } f$ and $f$ is continuous on $\text{int dom } f$
    \end{enumerate}
  \end{prop}

  \begin{prop}[Proposition 17.50]
    Let $f \in \Gamma_0(\mathcal{H})$ be Gâteaux differentiable at $x \in \text{dom} f$. Then $x \in \text{int dom } f$ and $f$ is continuous on $\text{int dom } f$.
  \end{prop}

  \begin{prop}[Proposition 17.51]
    Let $f: \mathcal{H} \to [-\infty, +\infty]$ be convex and suppose it is Gâteaux differentiable on some open subset $U$ of $\text{dom} f$. Then $f$ is continuous on $\text{int dom } f$.
  \end{prop}
\end{frame}

\begin{frame}
  \frametitle{Counterexample: Non-Continuous Gâteaux Differentiable Function}

  \begin{exmpl}[Example 17.49]
    Let $f$ be the discontinuous linear functional of Example 2.27, and set $g = |f|$. Then:
    \begin{itemize}
      \item $g$ is convex
      \item $g: \mathcal{H} \to \mathbb{R}$ is continuous at $0 = \text{dom} g$
      \item $g$ is Gâteaux differentiable at 0
      \item By Proposition 17.48, $g$ is lower semicontinuous at 0
      \item However, $g$ is neither continuous nor Fréchet differentiable at 0
    \end{itemize}
  \end{exmpl}

  This shows that Gâteaux differentiability alone does not guarantee continuity or Fréchet differentiability.
\end{frame}

\section{Python Numerical Examples}
\begin{frame}[fragile]
  \frametitle{Python: Computing Subdifferentials}

  \begin{lstlisting}
# Compute directional derivative numerically
def directional_derivative(f, x, y, eps=1e-8):
    alphas = np.logspace(-8, 0, 100)
    diffs = [(f(x + a*y) - f(x))/a
             for a in alphas]
    return np.min(diffs)

# Example: f(x) = x^2
f = lambda x: x**2
x0, y = 1.0, 2.0
f_prime = directional_derivative(f, x0, y)
print(f"f'({x0}; {y}) = {f_prime:.6f}")
  \end{lstlisting}
\end{frame}

\begin{frame}[fragile]
  \frametitle{Example: Convex Function Analysis}

  \begin{lstlisting}
class ConvexFunction:
    def __init__(self, f, grad_f):
        self.f = f
        self.grad_f = grad_f

    def is_monotone_gradient(self, x_vals):
        grads = [self.grad_f(x)
                 for x in x_vals]
        return all(grads[i] <= grads[i+1]
            for i in range(len(grads)-1))

# Example: f(x) = x^2 is convex
f1 = lambda x: x**2
grad_f1 = lambda x: 2*x
func1 = ConvexFunction(f1, grad_f1)
print(func1.is_monotone_gradient(
    np.linspace(-2, 2, 100)))
  \end{lstlisting}
\end{frame}

\section{Summary}
\begin{frame}
  \frametitle{Key Takeaways}

  \begin{enumerate}
    \item \textbf{Directional Derivatives} characterize behavior of convex functions at a point
    \item \textbf{Convexity Characterizations} provide first-order conditions via monotonicity of gradients
    \item \textbf{Gâteaux vs Fréchet Differentiability:}
      \begin{itemize}
        \item Gâteaux: Linear approximation in each direction
        \item Fréchet: Uniform linear approximation
        \item Equivalent in finite dimensions
      \end{itemize}
    \item \textbf{Subdifferentials} generalize gradients to non-smooth functions
    \item \textbf{Continuity:} Differentiability on interior implies continuity
    \item \textbf{Applications:} Optimization, monotone operators, descent methods
  \end{enumerate}
\end{frame}

\begin{frame}
  \frametitle{Further Reading}

  \begin{itemize}
    \item Proposition 17.2-17.4: Directional derivative properties
    \item Proposition 17.7-17.10: Convexity and strict convexity characterizations
    \item Proposition 17.21-17.22: Descent directions and steepest descent
    \item Proposition 17.31-17.41: Gâteaux and Fréchet differentiability
    \item Proposition 17.48-17.51: Continuity results
    \item Examples 17.8, 17.11, 17.13, 17.25, 17.33, 17.36, 17.46, 17.49: Important concrete examples
  \end{itemize}
\end{frame}

\end{document}
